% ============================================================
% UNIDAD 4: SISTEMAS INTEGRADOS G1000 Y G5000
% ============================================================
\clearpage
\section*{UNIDAD 4: SISTEMAS INTEGRADOS}
\subsection*{Garmin G1000 / G1000 NXi y la familia G3000 / G5000}

\begin{recuadro}
\textbf{Objetivo de la unidad:} conocer la arquitectura de un sistema de aviónica totalmente
integrado, sus unidades reemplazables, cómo se comunican, cómo se configuran y actualizan, y los
puntos clave de su mantenimiento.
\end{recuadro}

% ------------------------------------------------------------
\section*{1. EL GARMIN G1000}

Presentado en 2004, el G1000 fue el primer sistema totalmente integrado de uso masivo en la
aviación general. Equipa, entre otros, a los Cessna 172, 182, 206 y Caravan, Diamond DA40, DA42 y
DA62, Piper Seneca y Mooney, y versiones de turbohélices como el King Air. En 2017 apareció la
versión \textbf{G1000 NXi}, con procesadores y pantallas más rápidos y nuevas funciones, que también
se ofrece como actualización de los G1000 existentes.

\subsubsection*{1.1. Unidades del sistema (configuración típica de un avión monomotor)}
\begin{center}
\small
\begin{tabular}{|p{0.22\linewidth}|p{0.66\linewidth}|}
\hline
\textbf{Unidad} & \textbf{Función} \\ \hline
GDU (PFD y MFD) & Pantallas de 10,4 pulgadas (o 15 en otras instalaciones); idénticas entre sí,
la función la define su conexión. \\ \hline
GIA 63 / 63W (dos) & \textbf{Unidades de aviónica integrada}: cada una contiene un receptor GPS
(63W: con WAAS/SBAS), una radio VHF COM, un receptor VHF NAV y de senda de planeo, y funciona como
``centro de comunicaciones'' entre los demás equipos. \\ \hline
GDC 74 & Computadora de datos de aire. \\ \hline
GRS 77 / GRS 79 & AHRS. \\ \hline
GMU 44 & Magnetómetro, en el ala o la cola. \\ \hline
GEA 71 & Unidad de adquisición de datos del motor y la célula. \\ \hline
GMA 1347 & Panel de audio con intercomunicador y radiobaliza. \\ \hline
GTX 33 / GTX 345R & Transpondedor modo S remoto (con ADS-B en las versiones nuevas). \\ \hline
GFC 700 & Piloto automático integrado (servos en los comandos). \\ \hline
GDL & Receptor de datos meteorológicos por satélite (opcional). \\ \hline
\end{tabular}
\end{center}

\subsubsection*{1.2. Cómo se comunican}
\begin{center}
\begin{tikzpicture}[scale=0.92, transform shape, node distance=0.5cm, every node/.style={font=\footnotesize, align=center},
  b/.style={draw, thick, fill=ppAzul!8, minimum height=0.9cm, text width=1.9cm},
  s/.style={draw, thick, fill=orange!12, minimum height=0.8cm, text width=1.6cm}]
  \node[b] (pfd) at (0,3) {PFD (GDU)};
  \node[b] (mfd) at (6,3) {MFD (GDU)};
  \node[b] (g1) at (1.5,1.2) {GIA 1};
  \node[b] (g2) at (4.5,1.2) {GIA 2};
  \draw[very thick, ppAzul] (pfd) -- node[above] {Ethernet (HSDB)} (mfd);
  \draw[very thick, ppAzul] (pfd) -- (g1); \draw[very thick, ppAzul] (mfd) -- (g2);
  \draw[very thick, ppAzul] (pfd) -- (g2); \draw[very thick, ppAzul] (mfd) -- (g1);
  \node[s] (adc) at (-1.5,-0.8) {ADC};
  \node[s] (ahrs) at (0.5,-0.8) {AHRS + magnetóm.};
  \node[s] (gea) at (2.5,-0.8) {Motor (GEA)};
  \node[s] (xpdr) at (4.5,-0.8) {Transp.};
  \node[s] (aud) at (6.5,-0.8) {Audio};
  \node[s] (ap) at (8.4,-0.8) {Piloto autom.};
  \foreach \n in {adc,ahrs,gea} \draw[thick] (\n.north) -- (g1.south);
  \foreach \n in {gea,xpdr,aud,ap} \draw[thick] (\n.north) -- (g2.south);
  \draw[thick, dashed] (adc.north) -- (pfd.south west);
  \node at (3,-1.75) {\scriptsize Líneas finas: buses serie (RS-232, RS-485) y ARINC 429; línea
  gruesa: red Ethernet de alta velocidad (esquema simplificado)};
\end{tikzpicture}
\end{center}

Las pantallas y las GIA forman una red de alta velocidad. Los sensores y equipos periféricos se
conectan, generalmente de forma redundante, a las GIA y a las pantallas mediante buses serie y
ARINC 429. Si se pierde una GIA, el sistema sigue funcionando con la otra (con una radio y un GPS
menos).

% ------------------------------------------------------------
\section*{2. CONFIGURACIÓN, SOFTWARE Y BASES DE DATOS}

\begin{itemize}
  \item \textbf{Modo de configuración:} se ingresa encendiendo las pantallas con una tecla
        presionada (según el procedimiento del fabricante). Permite ver el estado de cada unidad,
        cargar software y configuración, calibrar sensores y leer registros de fallas.
  \item \textbf{Software y configuración:} se cargan desde una \textbf{tarjeta SD} preparada con la
        versión aprobada para esa aeronave. Cuando se reemplaza una unidad (LRU), hay que cargarle
        el software y la configuración correspondientes; si no coinciden, el sistema muestra
        alertas de incompatibilidad.
  \item \textbf{Bases de datos:} navegación (ciclo \textbf{AIRAC} de 28 días), terreno, obstáculos,
        planos de aeropuertos y cartas. Se actualizan por tarjeta SD o en forma inalámbrica en los
        sistemas nuevos. Volar IFR con una base de navegación vencida tiene limitaciones.
  \item \textbf{Registro:} toda carga de software o de configuración se registra en los registros
        de mantenimiento con la versión instalada.
\end{itemize}

% ------------------------------------------------------------
\section*{3. LA FAMILIA G3000 Y G5000}

Para aviones turbohélice y jets, Garmin desarrolló sistemas con pantallas más grandes y
\textbf{controladores táctiles}:
\begin{itemize}
  \item \textbf{G3000:} jets livianos y turbohélices (por ejemplo, HondaJet, Cessna Citation M2 y
        derivados en otros fabricantes).
  \item \textbf{G5000:} jets medianos, como los Cessna Citation Latitude, Longitude y Sovereign+ y
        los Learjet 70/75. Tres pantallas de gran tamaño (14 pulgadas) en formato apaisado, que se
        pueden dividir en ventanas, y \textbf{controladores con pantalla táctil} en el pedestal desde
        donde se manejan radios, plan de vuelo, sistemas y listas de chequeo.
  \item Integran funciones de aviones de transporte: piloto automático con acelerador automático en
        algunos modelos, sistema de alertas de la tripulación (CAS), páginas sinópticas de los
        sistemas (combustible, eléctrico, presurización) y planificación del rendimiento.
  \item Las versiones \textbf{Prime} recientes renovaron el software y la presentación, y aportaron
        funciones que luego se trasladaron a la aviación general.
\end{itemize}

% ------------------------------------------------------------
\section*{4. MANTENIMIENTO DE UN SISTEMA INTEGRADO}
\begin{itemize}
  \item Leer los \textbf{registros de fallas} y las páginas de estado antes de cambiar nada.
  \item Verificar primero \textbf{alimentación, conectores y cableado}: muchas fallas
        ``electrónicas'' son pines flojos, conectores mal asentados o masas deficientes.
  \item Al reemplazar una unidad: anotar número de parte, serie y versión de software; cargar
        software y configuración; realizar las calibraciones y ensayos que indica el manual (por
        ejemplo, AHRS y magnetómetro después de cambiar el AHRS o el magnetómetro).
  \item Ventilación: las pantallas tienen ventiladores; uno trabado provoca sobretemperatura y
        apagado.
  \item Actualizar las bases de datos y registrar la vigencia.
  \item Inspeccionar periódicamente las baterías de reserva y su prueba de capacidad.
\end{itemize}

% ------------------------------------------------------------
\section*{5. ACTIVIDADES DE LA UNIDAD}
\begin{enumerate}
  \item Dibujar el diagrama de interconexión de un G1000 de un Cessna 172 a partir del manual de
        mantenimiento del avión (capítulo 34 o 42, según el fabricante).
  \item ¿Qué funciones se pierden si falla la GIA 1? ¿Y si falla el GDC 74?
  \item Explicar el procedimiento para actualizar la base de datos de navegación y cómo se registra.
  \item Practicar con un entrenador o simulador de G1000 la carga de un plan de vuelo y la lectura
        de las páginas de motor.
  \item Comparar la cabina de un Cessna 172 con G1000 y la de una Citation Latitude con G5000:
        cantidad y tamaño de pantallas, forma de control, funciones.
\end{enumerate}
